This appendix contains supplementary empirical diagnostics referenced in Chapter~\ref{chap:results} but not required for the main sequence of results.

\section{Ridge Feature-Set Diagnostics}
\label{app:linear_feature_set_diagnostics}

Figure~\ref{fig:app-ridge-feature-deciles} compares validation errors across realized-drawdown deciles for the three Ridge feature sets. The three specifications have nearly identical error profiles. Errors are lowest in the middle of the realized-drawdown distribution and rise sharply in the most severe drawdown decile. The additional ownership and network features therefore do not materially change where the linear model succeeds or fails.

\begin{figure}[!htbp]
    \centering
    \includegraphics[width=0.82\textwidth]{04_01_ridge_validation_mae_by_realized_drawdown_decile.pdf}
    \caption[Ridge validation errors across realized drawdown deciles]{Validation MAE across realized-drawdown deciles for Ridge estimated using MKT, MKT~+~OWN, and MKT~+~OWN~+~NTW features.}
    \label{fig:app-ridge-feature-deciles}
\end{figure}

\section{XGBoost Feature-Set Diagnostics}
\label{app:nonlinear_feature_set_diagnostics}

Figure~\ref{fig:app-xgboost-feature-deciles} compares validation errors across realized-drawdown deciles for the three XGBoost feature sets. The nearly overlapping profiles show that adding ownership and network variables does not materially alter the distribution of prediction errors. Errors remain lowest in the middle of the realized-drawdown distribution and increase sharply for the most severe drawdowns.

\begin{figure}[!htbp]
    \centering
    \includegraphics[width=0.82\textwidth]{04_02_xgboost_validation_mae_by_realized_drawdown_decile.pdf}
    \caption[XGBoost validation errors across realized drawdown deciles]{Validation MAE across realized-drawdown deciles for XGBoost estimated using MKT, MKT~+~OWN, and MKT~+~OWN~+~NTW features.}
    \label{fig:app-xgboost-feature-deciles}
\end{figure}

\section{XGBoost Tuning Details}
\label{app:xgboost_tuning_details}

Table~\ref{tab:selected-xgboost-configuration} reports the complete configuration retained after the controlled tuning comparison in Section~\ref{sec:xgboost_tuning_results}.

\begin{table}[!htbp]
\centering
\footnotesize
\caption{Selected market-only XGBoost configuration.}
\label{tab:selected-xgboost-configuration}
\begin{tabular}{lr}
\toprule
Parameter & Value \\
\midrule
Learning rate & 0.10 \\
Maximum depth & 3 \\
Minimum child weight & 1 \\
Row subsample & 0.80 \\
Column subsample & 0.80 \\
Maximum boosting rounds & 1,000 \\
Early-stopping patience & 30 rounds \\
Training objective & Squared error \\
Validation metric & MAE \\
Selected boosting round & 90 \\
Training MAE & 0.0836 \\
Validation MAE & 0.0951 \\
\bottomrule
\end{tabular}
\end{table}


\section{GNN Architecture Details}
\label{app:gnn_architecture_details}

The three architectures are initially estimated using a common hidden dimension of 16, dropout of 0.20, learning rate of 0.001, and weight decay of 0.0001. Figure~\ref{fig:app-gnn-base-learning-curves} reports their individual training and validation learning curves.

\begin{figure}[!htbp]
    \centering
    \begin{subfigure}[t]{0.48\textwidth}
        \centering
        \includegraphics[width=\linewidth]{04_04_graphsage_learning_curve.pdf}
        \caption{GraphSAGE.}
    \end{subfigure}
    \hfill
    \begin{subfigure}[t]{0.48\textwidth}
        \centering
        \includegraphics[width=\linewidth]{04_04_gat_learning_curve.pdf}
        \caption{GAT.}
    \end{subfigure}

    \medskip

    \begin{subfigure}[t]{0.48\textwidth}
        \centering
        \includegraphics[width=\linewidth]{04_04_temporal_graphsage_learning_curve.pdf}
        \caption{Temporal GraphSAGE.}
    \end{subfigure}
    \caption{Learning curves for the three GNN architectures under the common initial configuration.}
    \label{fig:app-gnn-base-learning-curves}
\end{figure}

\section{Temporal GraphSAGE Tuning Details}
\label{app:gnn_tuning_details}

Table~\ref{tab:selected-temporal-graphsage-configuration} reports the complete configuration retained after the tuning comparison in Section~\ref{sec:temporal_graphsage_tuning}. Figure~\ref{fig:app-temporal-graphsage-tuning-curves} reports the individual learning curves for the four tuned variants.

\begin{table}[!htbp]
\centering
\footnotesize
\caption{Selected Temporal GraphSAGE configuration.}
\label{tab:selected-temporal-graphsage-configuration}
\begin{tabular}{lr}
\toprule
Parameter & Value \\
\midrule
Architecture & Temporal GraphSAGE \\
Hidden dimension & 32 \\
Learning rate & 0.001 \\
Dropout & 0.20 \\
Weight decay & 0.0001 \\
Optimizer & Adam \\
Training loss & MAE \\
Maximum epochs & 50 \\
Early-stopping patience & 5 epochs \\
Minimum improvement & 0.0001 \\
Selected epoch & 28 \\
Training MAE & 0.0813 \\
Validation MAE & 0.0941 \\
\bottomrule
\end{tabular}
\end{table}


\begin{figure}[!htbp]
    \centering
    \begin{subfigure}[t]{0.48\textwidth}
        \centering
        \includegraphics[width=\linewidth]{04_04_temporal_graphsage_hidden_16_lr_0010_learning_curve.pdf}
        \caption{Hidden 16; learning rate 0.010; dropout 0.20.}
    \end{subfigure}
    \hfill
    \begin{subfigure}[t]{0.48\textwidth}
        \centering
        \includegraphics[width=\linewidth]{04_04_temporal_graphsage_hidden_32_learning_curve.pdf}
        \caption{Hidden 32; learning rate 0.001; dropout 0.20.}
    \end{subfigure}

    \medskip

    \begin{subfigure}[t]{0.48\textwidth}
        \centering
        \includegraphics[width=\linewidth]{04_04_temporal_graphsage_hidden_32_dropout_030_learning_curve.pdf}
        \caption{Hidden 32; learning rate 0.001; dropout 0.30.}
    \end{subfigure}
    \hfill
    \begin{subfigure}[t]{0.48\textwidth}
        \centering
        \includegraphics[width=\linewidth]{04_04_temporal_graphsage_hidden_64_dropout_030_learning_curve.pdf}
        \caption{Hidden 64; learning rate 0.001; dropout 0.30.}
    \end{subfigure}
    \caption{Learning curves for the four tuned Temporal GraphSAGE configurations.}
    \label{fig:app-temporal-graphsage-tuning-curves}
\end{figure}

\section{XGBoost Interpretation Diagnostics}
\label{app:model_interpretation_details}

Figure~\ref{fig:app-xgboost-shap-dependence} reports SHAP dependence plots for the three most influential predictors. Higher downside volatility and larger past maximum drawdown are associated with larger positive contributions to predicted downside risk, whereas larger market capitalization generally contributes negatively. These patterns describe the fitted model and should not be interpreted as causal effects.

\begin{figure}[!htbp]
    \centering
    \begin{subfigure}[t]{0.48\textwidth}
        \centering
        \includegraphics[width=\linewidth]{04_11_xgboost_shap_dependence_1_downside_volatility_63_days.pdf}
        \caption{Downside volatility.}
    \end{subfigure}
    \hfill
    \begin{subfigure}[t]{0.48\textwidth}
        \centering
        \includegraphics[width=\linewidth]{04_11_xgboost_shap_dependence_2_maximum_drawdown_252_days.pdf}
        \caption{Past maximum drawdown.}
    \end{subfigure}

    \medskip

    \begin{subfigure}[t]{0.48\textwidth}
        \centering
        \includegraphics[width=\linewidth]{04_11_xgboost_shap_dependence_3_log_market_capitalization.pdf}
        \caption{Market capitalization.}
    \end{subfigure}
    \caption[SHAP dependence patterns for the leading predictors]{SHAP contributions across the transformed values of the three most influential predictors in the selected XGBoost model. The fitted lines summarize the corresponding dependence patterns.}
    \label{fig:app-xgboost-shap-dependence}
\end{figure}

\noindent
Figure~\ref{fig:app-xgboost-quarterly-shap} compares normalized SHAP importance across test quarters. The feature ordering is highly stable over time, showing that the selected model relies on a consistent set of predictors rather than changing its principal drivers from one quarter to another.

\begin{figure}[!htbp]
    \centering
    \includegraphics[width=0.9\textwidth]{04_11_xgboost_quarterly_shap_importance_heatmap.pdf}
    \caption[Feature-importance stability across test quarters]{Normalized mean absolute SHAP importance by test quarter for the selected XGBoost model. Each column sums to one, and features are ordered by their average importance.}
    \label{fig:app-xgboost-quarterly-shap}
\end{figure}

\section{Held-Out Model Comparison Diagnostics}
\label{app:quarterly_test_performance}

Figure~\ref{fig:finalist-test-deciles} compares the three finalists across realized-drawdown deciles. Their error profiles remain close through the first nine deciles and rise sharply in the tenth, showing that all three models have greater difficulty estimating the magnitude of the most severe downside outcomes.

\begin{figure}[!htbp]
    \centering
    \includegraphics[width=0.82\textwidth]{04_08_test_mae_by_realized_drawdown_decile.pdf}
    \caption[Held-out errors across realized drawdown deciles]{Test MAE across realized-drawdown deciles for the three frozen finalist models.}
    \label{fig:finalist-test-deciles}
\end{figure}

\noindent
Figure~\ref{fig:app-finalist-quarterly-performance} reports the MAE and within-quarter Spearman correlation of the three frozen finalist models in each test quarter. Quarterly MAE varies across the test period, while within-quarter rank correlations remain consistently positive. This indicates that the models' ability to rank stocks is more stable than their ability to estimate the magnitude of realized downside across periods.

\begin{figure}[!htbp]
    \centering
    \begin{subfigure}[t]{0.48\textwidth}
        \centering
        \includegraphics[width=\linewidth]{04_08_quarterly_test_mae_finalists.pdf}
        \caption{Quarterly MAE.}
        \label{fig:app-finalist-quarterly-mae}
    \end{subfigure}
    \hfill
    \begin{subfigure}[t]{0.48\textwidth}
        \centering
        \includegraphics[width=\linewidth]{04_08_quarterly_test_spearman_finalists.pdf}
        \caption{Quarterly rank correlation.}
        \label{fig:app-finalist-quarterly-spearman}
    \end{subfigure}
    \caption[Quarterly held-out performance of the finalist models]{MAE and Spearman rank correlation by test quarter for the three frozen finalist models.}
    \label{fig:app-finalist-quarterly-performance}
\end{figure}

\section{Firm-Size Error Diagnostic}
\label{app:firm_size_errors}

Figure~\ref{fig:app-xgboost-mae-size} reports the selected model's held-out MAE across within-quarter market-capitalization quintiles. Errors are larger among smaller stocks, although this supplementary relationship is not subjected to a separate robustness analysis.

\begin{figure}[!htbp]
    \centering
    \includegraphics[width=0.72\textwidth]{04_09_xgboost_test_mae_by_size.pdf}
    \caption[Selected-model errors by firm size]{Held-out MAE across within-quarter market-capitalization quintiles.}
    \label{fig:app-xgboost-mae-size}
\end{figure}

\section{Predictive Robustness Diagnostics}
\label{app:predictive_robustness_stability}

Figure~\ref{fig:app-predictive-robustness-deciles} reports standardized realized outcomes across predicted-risk deciles. The monotonic or near-monotonic profiles show that the model's ranking ability extends across downside-risk definitions and horizons.

\begin{figure}[!htbp]
    \centering
    \includegraphics[width=0.88\textwidth]{04_10_xgboost_robustness_decile_profiles.pdf}
    \caption[Predicted-risk decile profiles across robustness targets]{Standardized realized outcomes across predicted-risk deciles for the primary and alternative targets.}
    \label{fig:app-predictive-robustness-deciles}
\end{figure}

\noindent
Figure~\ref{fig:app-predictive-robustness-quarterly-spearman} reports quarterly ranking performance for the primary and alternative targets. Figure~\ref{fig:app-predictive-robustness-quarterly-mae} reports each quarter's MAE relative to the corresponding target's overall test MAE.

\begin{figure}[!htbp]
    \centering
    \includegraphics[width=0.82\textwidth]{04_10_xgboost_robustness_quarterly_spearman.pdf}
    \caption[Quarterly ranking performance across robustness targets]{Quarterly Spearman correlations between predicted and realized risk for the primary and alternative targets. Each point represents one test quarter, and diamonds mark the corresponding target means.}
    \label{fig:app-predictive-robustness-quarterly-spearman}
\end{figure}

\begin{figure}[!htbp]
    \centering
    \includegraphics[width=0.82\textwidth]{04_10_xgboost_robustness_quarterly_relative_mae.pdf}
    \caption[Quarterly error stability across robustness targets]{Quarterly MAE relative to each target's overall test MAE. The dashed line at one denotes the target-specific average.}
    \label{fig:app-predictive-robustness-quarterly-mae}
\end{figure}

\section{Outcomes by Fragility Score Decile}
\label{app:fragility_score_supporting_outcomes}

Figure~\ref{fig:fragility-score-deciles} shows that predicted and realized maximum drawdown rise across Fragility Score deciles, with the strongest separation in the upper tail.

\begin{figure}[!htbp]
    \centering
    \includegraphics[width=0.78\textwidth]{04_11_fragility_score_deciles.pdf}
    \caption[Predicted and realized drawdown across Fragility Score deciles]{Mean predicted and realized 63-day maximum drawdown by within-quarter Fragility Score decile in the held-out sample.}
    \label{fig:fragility-score-deciles}
\end{figure}

\noindent
Figure~\ref{fig:app-fragility-score-supporting-outcomes} shows that higher Fragility Score deciles also exhibit greater realized downside volatility and worse five-day losses, while cumulative returns weaken toward the most fragile decile. The score therefore captures a broader downside-risk gradient rather than mechanically ranking only the primary maximum-drawdown target.

\begin{figure}[!htbp]
    \centering
    \includegraphics[width=\textwidth]{04_11_fragility_score_supporting_outcomes.pdf}
    \caption[Supporting outcomes across Fragility Score deciles]{Mean realized downside volatility, worst five-day loss, and cumulative return by within-quarter Fragility Score decile.}
    \label{fig:app-fragility-score-supporting-outcomes}
\end{figure}

\section{Primary Portfolio Factor Regressions}
\label{app:primary_portfolio_alpha}

Table~\ref{tab:portfolio-alpha} reports the complete FF5-plus-momentum regression statistics underlying Figure~\ref{fig:portfolio-strategy-alpha}.

\begin{table}[!htbp]
\centering
\footnotesize
\caption{FF5-plus-momentum alpha for the primary portfolio comparison.}
\label{tab:portfolio-alpha}
\begin{tabular}{lrrrr}
\toprule
Portfolio & Annualized alpha & $t$-statistic & $p$-value & Adjusted $R^2$ \\
\midrule
All eligible stocks & -5.2\% & -1.35 & 0.179 & 0.946 \\
Exclude Decile 10 & 0.0\% & 0.01 & 0.992 & 0.974 \\
Screened minus all stocks & 5.2\% & 2.77 & 0.006 & 0.174 \\
\bottomrule
\end{tabular}
\end{table}


\section{Quarterly Portfolio Risk Comparison}
\label{app:portfolio_window_risk}

Figure~\ref{fig:portfolio-event-window-risk} compares the average downside outcomes of the four primary portfolios across the quarterly holding windows. The screened portfolio improves modestly on the full equal-weighted universe, while the much larger separation between Decile~1 and Decile~10 confirms the score's economic ordering power.

\begin{figure}[!htbp]
    \centering
    \includegraphics[width=\textwidth]{04_11_fragility_portfolio_risk.pdf}
    \caption[Risk across Fragility Score portfolios]{Average maximum drawdown, downside volatility, and worst five-day loss across the quarterly portfolio windows.}
    \label{fig:portfolio-event-window-risk}
\end{figure}

\section{Long-Only Portfolio Robustness}
\label{app:portfolio_robustness_figure}

Figure~\ref{fig:portfolio-robustness} compares the full and screened portfolios under alternative exclusion thresholds and weighting rules. The equal-weighted results remain similar when the exclusion threshold is widened, whereas the value-weighted portfolios are nearly unchanged by the screen.

\begin{figure}[!htbp]
    \centering
    \includegraphics[width=\textwidth]{04_12_portfolio_robustness.pdf}
    \caption[Long-only portfolio robustness]{Performance of the full and screened portfolios under alternative exclusion thresholds and weighting rules.}
    \label{fig:portfolio-robustness}
\end{figure}

\section{Long--Short Transaction-Cost Sensitivity}
\label{app:long_short_cost_sensitivity}

Table~\ref{tab:long-short-cost-sensitivity} reports annualized returns under alternative one-way transaction-cost assumptions. The volatility-managed calculation includes turnover from both quarterly DFRAG-10 reconstitution and changes in volatility-targeted exposure. The range illustrates sensitivity to implementation costs rather than claiming that any single rate precisely represents realized trading costs.

\begin{table}[!htbp]
\centering
\footnotesize
\caption{Annualized returns of the long--short strategies under alternative one-way transaction-cost assumptions.}
\label{tab:long-short-cost-sensitivity}
\begin{tabular}{lrrrr}
\toprule
Strategy & Gross & 10 bps & 25 bps & 50 bps \\
\midrule
DFRAG-10 & 55.7\% & 55.1\% & 54.2\% & 52.7\% \\
DFRAG-20 & 16.0\% & 15.6\% & 15.0\% & 14.0\% \\
DFRAG-10: volatility managed & 28.2\% & 27.5\% & 26.5\% & 24.9\% \\
\bottomrule
\end{tabular}
\end{table}

